\section{Closed transport tests and selection of a branch}
\label{test:section}

We translate a future $C^1$ exit into a transport bound shared by an
open set of finite-duration geodesics. Each such bound defines a closed
set of data, indexed by a basic open seed set and one integer. These
tests concern the full development and are independent of the Kerr
parameters. We then show that every passing test supplies a geodesic
approaching a noncorner branch, where the later perturbation will be
placed. For this selection we combine the comparison development of
\cite[\CompanionRef{foundation}{sec:comparison}]{Foundation} with
conservation of canonical phase-space volume. No curvature bound along
the selected observer is needed.

\subsection{Seeds, positive norms, and the passing sets}

A seed $w\in T\Sigma$ includes its base point. For $d=(h,K)$ let
$n$ be the future unit normal to the initial hypersurface in
$\mathcal M_d$, and define
\begin{equation}\label{test:seed}
 U_d(w)=\sqrt{1+h(w,w)}\,n+w,
 \qquad
 \gamma_{d,w}:[0,T_d(w))\longrightarrow\mathcal M_d .
\end{equation}
Here $\gamma_{d,w}$ is the maximal future unit timelike geodesic with
this initial velocity, and $T_d(w)\in(0,\infty]$ is its maximal
proper duration in the \emph{full} maximal globally hyperbolic
development. On $\mathbb R n\oplus T\Sigma$ put
\begin{equation}\label{test:positive-metric}
 k_{d,w}(0)(sn+X,tn+Y)=st+h(X,Y),
 \qquad \nabla_{\dot\gamma_{d,w}}k_{d,w}=0.
\end{equation}
Thus $k_{d,w}(\tau)$ is positive definite even though the spacetime
metric is Lorentzian. All operator norms below are induced by this
positive inner product.

For a piecewise smooth loop $\ell:[0,1]\to\mathcal M_d$ based at
$p=\gamma_{d,w}(\tau)$, parallel transport $k_{d,w}(\tau)$ along
$\ell$ and call the resulting positive inner products $k_\ell(s)$.
Define the developed length and based transport operator by
\begin{equation}\label{test:developed-length}
 L(\ell)=\int_0^1|\dot\ell(s)|_{k_\ell(s)}\,ds,
 \qquad P_\ell:T_p\mathcal M_d\longrightarrow T_p\mathcal M_d.
\end{equation}
If $F(s)$ is the parallel frame along $\ell$ starting from a
$k_{d,w}(\tau)$-orthonormal frame, and $\dot\ell=Fz$, then
$L(\ell)=\int_0^1|z(s)|\,ds$. The metric at the end of the
loop need not equal its initial value; Equation~\eqref{test:developed-length}
uses the transported metric at each parameter value.

\begin{lemma}[Comparison with the observer's positive metric]
\label{test:observer-comparison}
Let $U=\dot\gamma_{d,w}$ and set
$r_U=g_d+2U^\flat\otimes U^\flat$. This positive metric is parallel
along $\gamma_{d,w}$ and satisfies
\begin{equation}\label{test:observer-norm-bound}
 c(w)^{-1}k_{d,w}\le r_U\le c(w)k_{d,w},
 \qquad
 c(w)=\bigl(\sqrt{1+h(w,w)}+\sqrt{h(w,w)}\bigr)^2.
\end{equation}
In particular the comparison is uniform when the initial data and
seeds range in a sufficiently small neighborhood of a fixed pair.
\end{lemma}

\begin{proof}
Both $g_d$ and $U$ are parallel, so $r_U$ is parallel. At the initial
point choose an $h$-orthonormal spatial frame with first vector along
$w$, if $w\ne0$, and write $s=h(w,w)$ and
$\alpha=\sqrt{1+s}$. In the normal and first spatial directions,
the matrix of $r_U$ relative to $k_{d,w}(0)$ is
\[
 \begin{pmatrix}1+2s&-2\alpha\sqrt{s}\\
                 -2\alpha\sqrt{s}&1+2s\end{pmatrix}.
\]
Its eigenvalues are $(\alpha+\sqrt{s})^2$ and
$(\alpha-\sqrt{s})^2$, which are reciprocal. The other two
eigenvalues are one. This also covers $s=0$, and parallel transport
preserves the inequalities.
\end{proof}

Fix a countable basis $\mathcal O$ of nonempty relatively compact
open subsets of the six-dimensional manifold $T\Sigma$.
The test below asks for one transport constant at every positive time
on every geodesic in a basis set. In an extension this uniformity will
come from a compact earlier geodesic tube together with one regular
chart near the exit.

\begin{definition}[Passing test]\label{test:definition}
For $O\in\mathcal O$ and $B\in\mathbb N$, $B\ge1$, the set
$\mathcal A(O,B)\subset\mathcal U$ consists of the data such that,
for every $w\in O$, one has $T_d(w)\le B$ and, for every
$0<\tau<T_d(w)$ and every piecewise smooth loop $\ell$ based at
$\gamma_{d,w}(\tau)$ whose entire image is contained in
$I^+_{\mathcal M_d}(\Sigma)$,
\begin{equation}\label{test:inequality}
 L(\ell)<B^{-1}
 \quad\Longrightarrow\quad
 \lVert P_\ell-\operatorname{Id}\rVert_{k_{d,w}(\tau)}
       \le B L(\ell).
\end{equation}
\end{definition}

\begin{lemma}[Closedness]\label{test:closed}
Each $\mathcal A(O,B)$ is relatively closed in $\mathcal U$.
\end{lemma}

\begin{proof}
We show that every failure persists in a neighborhood of the datum.
If $T_d(w)>B$ for one $w\in O$, choose
$B<b<T_d(w)$. The geodesic segment $\gamma_{d,w}([0,b])$ is a
compact witness. Otherwise, failure is witnessed by one
$w\in O$, one $0<\tau<T_d(w)$, and one loop with the two strict
inequalities
\[
 L(\ell)<B^{-1},\qquad
 \lVert P_\ell-\operatorname{Id}\rVert_{k_{d,w}(\tau)}>B L(\ell).
\]

Smooth compact-subset vacuum Cauchy stability and geometric uniqueness
apply to these witnesses: the data are smooth positive vacuum
constraint data, and convergence in all $p_m$ implies smooth
convergence on every compact part of $\Sigma$. This is precisely the
compact-subset Cauchy statement used in
\cite[proof of \CompanionRef{foundation}{cat:closed-tests}]{Foundation}. To specify its use, enlarge the
compact witness by finitely many local evolution neighborhoods and
their causal shadows on the Cauchy hypersurface. The shadows of a
compact subset of a globally hyperbolic spacetime on a Cauchy
hypersurface are compact. Finite local Cauchy stability, followed by
geometric uniqueness on the overlaps, therefore gives embeddings of
a neighborhood of the witness into the nearby developments on which
the metrics converge smoothly. Only finitely many finite orders are
needed for any fixed witness.

For a loop witness also include compact timelike arcs from $\Sigma$
to finitely many neighborhoods covering its image. Such neighborhoods
exist because the original image is compact and contained in
$I^+(\Sigma)$. After a small change of data those arcs remain timelike;
their final short segments can be varied within the neighborhoods.
Consequently the nearby loop images are still contained in
$I^+(\Sigma)$.

The initial velocity in Equation~\eqref{test:seed}, the geodesic on
the chosen compact time interval, and parallel transport all vary
continuously, by smooth ODE dependence in these common neighborhoods.
For the loop, use a local diffeomorphism converging smoothly to the
identity to move its base to the nearby geodesic at the same proper
time; its support can be chosen in one of the preceding neighborhoods.
Lengths and transport matrices then converge, so both strict
inequalities persist. The duration witness similarly persists beyond
$B$. Thus the complement is open. No continuity assertion for
$T_d(w)$ at a maximal endpoint is used.
\end{proof}

\subsection{Covering future exits}

The next three lemmas convert a single extension exit into a uniform
test on an open family of initial seeds. Causal homotopy forces these
geodesics to have finite duration in the full development, while a
small-developed-length bootstrap supplies the common transport bound.

\begin{lemma}[Intrinsic diamond under a causal homotopy]
\label{test:diamond}
Let $P$ be a smooth globally hyperbolic open subspacetime of a
Hausdorff manifold $\widetilde M$ carrying a continuous Lorentzian
metric. Let $H:[0,1]^2\to\widetilde M$ be continuous, with
$H(\lambda,0)=p$ and $H(\lambda,1)=q$ for fixed $p,q\in P$.
Suppose every $H(\lambda,\cdot)$ is future causal. If the path for
$\lambda=0$ lies in $P$, then all the paths lie in $P$.
\end{lemma}

\begin{proof}
Let $E$ be the set of parameters whose whole paths are in $P$.
Compactness of $[0,1]$ and openness of $P$ imply that $E$ is open.
For $\lambda\in E$ the path lies in the intrinsic diamond
$D=J_P^+(p)\cap J_P^-(q)$, which is compact by global
hyperbolicity. Its image in $\widetilde M$ is compact and hence
closed because the ambient manifold is Hausdorff. If
$\lambda_j\in E$ and $\lambda_j\to\lambda$, continuity of $H$
puts every $H(\lambda,s)$ in that same closed image of $D$, so
$\lambda\in E$. Thus $E$ is nonempty, open, and closed in the
connected interval. This argument does not assume ambient causal
convexity or any ambient geodesic uniqueness.
\end{proof}

\begin{lemma}[Interior endpoints of timelike geodesics]
\label{test:interior-endpoint}
A maximal timelike geodesic in a smooth spacetime has no endpoint
there as an unparametrized causal curve. Its maximal past continuation
from a point in $I^+(\Sigma)$ therefore crosses a Cauchy
hypersurface $\Sigma$.
\end{lemma}

\begin{proof}
Suppose a unit geodesic has a position limit at an endpoint in the
spacetime. Choose a precompact temporal coordinate chart $(t,x)$
about the limit. On a sufficiently late tail, the causal slopes
$V=(1,dx/dt)$ are uniformly bounded. With $\dot t=dt/d\tau>0$,
the geodesic equation reads
\begin{equation}\label{test:time-geodesic}
 \frac{d}{dt}\log\dot t
   =-\Gamma^0_{\alpha\beta}V^\alpha V^\beta,
 \qquad
 \frac{dV^i}{dt}
   =-\Gamma^i_{\alpha\beta}V^\alpha V^\beta
       +\Gamma^0_{\alpha\beta}V^\alpha V^\beta V^i.
\end{equation}
The right sides are bounded on the finite chart-time tail. Hence
$\dot t$ is bounded above and below by positive constants and has
a limit; $V$ also has a limit. The remaining proper time is finite,
and the limiting position and unit timelike tangent give smooth ODE
continuation. This contradicts maximality. The same proof applies at
the past end. Complete a maximal past geodesic by any future
inextendible timelike continuation. Cauchyness gives a crossing of
$\Sigma$, and achronality places that crossing to the past of a
point in $I^+(\Sigma)$.
\end{proof}

\begin{lemma}[Arbitrary loops of small developed length]
\label{test:loop-bootstrap}
Consider a family of base points in regular coordinate charts for a
$C^1$ Lorentzian metric. Suppose the Christoffel symbols are bounded
by one constant, each base point has coordinate distance at least
$\rho>0$ from its chart boundary, and the coordinate matrices of
the prescribed positive orthonormal frames and their inverses have
norm at most $K$. There are $l_*>0$ and $C_*<\infty$, depending
only on these constants, such that every piecewise smooth loop based
at any of these points with developed length $L<l_*$ remains in its
selected chart and satisfies
\begin{equation}\label{test:local-loop-bound}
 \lVert P_\ell-\operatorname{Id}\rVert_k\le C_*L.
\end{equation}
The loop may be specified in any open subspacetime of the chart's
ambient manifold; no initial assumption of coordinate smallness is
required.
\end{lemma}

\begin{proof}
This is the bounded-connection parallel-transport argument underlying
\cite[Lemma~2.13]{Sbierski2022}, with a first-exit bootstrap to handle
a loop specified only by its developed length. Related comparisons
with coordinate length, including the small-developed-length case,
appear in \cite[Lemma~4.2 and Corollary~4.12]{Sbierski2024}.
Continue the initial frame $F_0$ parallel along the loop, and write
$\dot\ell=Fz$. Up to the first departure from the selected chart,
the coordinate parallel-transport equation implies, for a fixed
constant $C$ including the dimension-dependent matrix factors,
\[
 |\dot F|\le C\lVert F\rVert^2|z|,
 \qquad |\dot\ell|\le\lVert F\rVert|z|.
\]
Use accumulated developed length
$s=\int|z|$ as the parameter on nonconstant pieces, or equivalently
integrate these inequalities without reparametrizing. Comparison with
$f'=Cf^2$, $f(0)=K$, gives
$\lVert F\rVert\le K/(1-CKs)\le2K$ for
$s\le(2CK)^{-1}$. Thus coordinate displacement is at most $2Ks$.
Choose
\[
 l_*<\min\{(2CK)^{-1},\rho/(4K)\}.
\]
A first chart exit at developed length less than $l_*$ would have
displacement at least $\rho$ and at most $\rho/2$, a
contradiction. The estimates consequently hold on the entire loop.
They give
$\lVert F(1)-F_0\rVert\le4CK^2L$. Since the loop returns to its
base point, its transport matrix in the initial orthonormal frame is
$F_0^{-1}F(1)$. Multiplication by the uniformly bounded inverse proves
Equation~\eqref{test:local-loop-bound}, for example with
$C_*=4CK^3$. Piecewise smooth parametrizations cause no change in
these integral estimates.
\end{proof}

\begin{proposition}[Coverage of future $C^1$ exits]
\label{test:exit-cover}
If $\mathcal M_d$, $d\in\mathcal U$, admits a future $C^1$
extension in the stated sense, then $d\in\mathcal A(O,B)$ for
some $O\in\mathcal O$ and some integer $B\ge1$.
\end{proposition}

\begin{proof}
Identify $\mathcal M_d$ with its open image in the extension.
Let $c$ be the timelike curve reaching its boundary. Its tail has no
interior future endpoint, since an interior limit and the given
ambient boundary limit would coincide in the Hausdorff extension.
Extending the interior part maximally to the past and using the
Cauchy property, we find that a sufficiently late part lies in
$I^+_{\mathcal M_d}(\Sigma)$.

\emph{A uniformly timelike straight segment reaching the boundary.}
We first obtain a straight timelike segment that stays inside
$\mathcal M_d$ until its final point. The construction also covers
an exiting curve whose timelike margin tends to zero at its endpoint.
Take a precompact ambient chart $(t,x)$ containing the endpoint and
a tail of $c$, with $dt$ temporal and $\partial_t$ future timelike.
In chart time the tail is $(t,x(t))$, $t\le t_f$. On a smaller
compact coordinate box, causal slopes $v=dx/dt$ are bounded, and
there is $c_0>0$ such that
$\widetilde g((1,v),\partial_t)\le-c_0$ for all these slopes.
The metric is Lipschitz there. Replace the curve's time coordinate by
\[
 \widetilde t(t)=t-\epsilon_1(e^{H(t_f-t)}-1),
 \qquad \epsilon_1>0.
\]
Put $\Delta=t-\widetilde t$ and
$\eta=\widetilde t'-1=\epsilon_1H e^{H(t_f-t)}$.
Then $0\le\Delta\le\eta/H$. The change of the quadratic form
from the displacement is bounded by $C\Delta(1+\eta)^2$, whereas
the addition of $\eta\partial_t$ contributes at most
$-2c_0\eta$; its quadratic contribution is nonpositive.
Choose $H$ large and then $\epsilon_1$ small enough that
$\eta\le1$. The new velocity is uniformly future timelike, with
quadratic form at most $-c_0\epsilon_1H$ after increasing $H$ if
necessary. The new curve has the same final point and begins inside
$I^+(\Sigma)$, since its initial displacement can be made
arbitrarily small.

After reparametrization by its new chart time, take a sufficiently
fine polygonal interpolation with the same endpoints. It remains
uniformly future timelike. Indeed, on a short subinterval the
velocities lie in a bounded convex cone of future timelike vectors
with a fixed smaller margin for one frozen metric; averaging preserves
that cone, and continuity of the metric preserves a smaller margin
on the interpolating segment. A finite subdivision suffices.
Follow this polygon to its first exit from $\mathcal M_d$ and let
$z$ be that exit. The straight piece immediately before $z$ has
constant spatial slope $v_*$, is interior before $z$, and is uniformly
timelike up to $z$.

\emph{An open family of intrinsic geodesics has finite full duration.}
Choose $l>0$ small and set $t_0=t(z)-l$, $t_1=t(z)+l$ so that
the last length-$l$ part of the straight piece is in the chart.
Start intrinsic unit timelike geodesics at $t=t_0$ with position
within $l^2$ of that piece and spatial slope within $l$ of $v_*$,
restricting to a nonempty open subset of states inside
$I^+_{\mathcal M_d}(\Sigma)$. All these states form a
six-dimensional open family. The Christoffel symbols of the ambient
$C^1$ metric are continuous and bounded. While an intrinsic geodesic
remains in the box, Equation~\eqref{test:time-geodesic} therefore
bounds its spatial acceleration by a constant $C_0$, uniformly over
the family. On the interval of chart-time length $2l$,
\begin{equation}\label{test:near-straight}
 |V^i-v_*^i|\le (1+2C_0)l,
 \qquad
 |x(t)-x_*(t)|\le C_1l^2,
\end{equation}
where $x_*$ is the straight line extended through the box. Choose a
compact box about that line with a fixed positive margin and shrink
$l$ first. These estimates, by a first-exit argument, prevent chart
departure before $t_1$. They also keep all slopes in a common
uniformly timelike neighborhood of $v_*$. Unit normalization gives
uniform positive upper and lower bounds for $dt/d\tau$.

If one such geodesic remains in $\mathcal M_d$ through $t_1$,
Equation~\eqref{test:near-straight} places its point at $t=t(z)$
within $C_1l^2$ of $z$. Add to its spatial path a smooth bump
supported in $(t_0,t_1)$ that moves this point to $z$ and leaves the
endpoints fixed. A rescaled fixed bump has height $O(l^2)$ and
derivative $O(l)$. Multiplying the bump by
$\lambda\in[0,1]$ gives a fixed-endpoint homotopy through future
timelike curves inside the ambient box. Lemma~\ref{test:diamond},
applied to the globally hyperbolic spacetime $\mathcal M_d$,
forces the final path to remain inside it, although it passes through
$z\notin\mathcal M_d$. This is impossible.

Thus every member of the family has maximal intrinsic chart-time
strictly less than, or equal to, $t_1$. The preceding bounds imply
finite remaining proper time and limits of position and unit tangent
in the ambient box. The position limit lies outside $\mathcal M_d$:
otherwise the smooth intrinsic geodesic equation would continue it.
This is its maximal endpoint in the \emph{full} $\mathcal M_d$.
Any intrinsic continuation would have an interior position whose
image is the already determined ambient boundary limit, a
contradiction. No geodesic equation has been solved outside
$\mathcal M_d$, and no uniqueness assertion for the merely
continuous ambient Christoffel symbols has been made.

Continue one of the starting states backwards in $\mathcal M_d$.
Lemma~\ref{test:interior-endpoint} and its location in $I^+(\Sigma)$
give a transverse crossing of $\Sigma$. Smooth geodesic flow through
this fixed compact earlier segment maps an open neighborhood in
$T\Sigma$ diffeomorphically onto an open part of the above
starting-state family. Shrink it to a relatively compact open set
with closure mapped into that family. Its earlier geodesic tube has
compact closure and uniformly bounded proper duration. The final
chart interval has uniformly bounded proper duration as well. Hence
there is $B_0<\infty$ such that $T_d(w)\le B_0$ for every seed
in this smaller neighborhood.

\emph{One transport constant for all base points and all short loops.}
Initial frames orthonormal for Equation~\eqref{test:positive-metric}
and their inverses are uniformly bounded on the compact seed closure.
On the compact earlier tube, a finite cover by regular charts and
the parallel-transport ODE give uniform bounds for their continued
matrices and inverses. On the final part use the ambient chart:
bounded Christoffel symbols and bounded chart-time slopes on an
interval of length $2l$ give the same bounds by the linear
parallel-transport equation and its inverse equation. All final base
points stay in a compact coordinate box strictly inside this chart,
even as they approach the boundary of the image of $\mathcal M_d$.
A finite chart cover of the earlier compact tube, with smaller chart
cores, provides a common positive chart margin there as well.
Lemma~\ref{test:loop-bootstrap} now supplies one $l_*>0$ and one
$C_*<\infty$ for every base point on every geodesic in the seed
neighborhood. In particular it applies to every loop required in
Definition~\ref{test:definition}.

Choose an integer $B\ge\max\{1,B_0,C_*,l_*^{-1}\}$ and a
basis set $O\in\mathcal O$ inside this neighborhood. The strict
premise $L<B^{-1}$ implies $L<l_*$, so all the requirements for
$d\in\mathcal A(O,B)$ hold.
\end{proof}

\subsection{Comparison runout and the corner estimates}

The next input concerns the fixed arbitrary subextremal center. Its
smallness threshold is chosen once, before the seed set, the passing
constant, or any later perturbation order.

\begin{proposition}[Comparison development and first runout]
\label{test:comparison}
After decreasing the defining $p_{10}$ radius of $\mathcal U$, each
$d\in\mathcal U$ has a smooth globally hyperbolic vacuum comparison
development $\mathcal P_d$ containing the whole bridge $\Sigma$ as
a Cauchy hypersurface and embedded openly and isometrically in
$\mathcal M_d$. It contains a spacelike cylinder $\mathcal S$ and
its full future product wedge, with coordinates
$t=(u+v)/2\ge0$, $x=(v-u)/2\in\mathbb R$,
$\theta\in\mathbb S^2$, in which
\begin{equation}\label{test:wedge}
 \begin{split}
 g&=-2a\,du\,dv+
 \gamma_{AB}(d\theta^A-b^Adu)(d\theta^B-b^Bdu),\\
 a&=q\mathcal A,\qquad
 q=q_0e^{-\kappa(u+v)},\qquad q_0,\kappa>0.
 \end{split}
\end{equation}
The coefficients $\mathcal A^{\pm1}$ and $\gamma^{\pm1}$ are
uniformly bounded, with $\gamma$ uniformly comparable to a fixed
round sphere metric. Every seed with $T_d(w)<\infty$ crosses
$\mathcal S$ before its first runout from $\mathcal P_d$, and
that first runout has exactly one of the following coordinate limits:
\begin{equation}\label{test:runout}
 \begin{array}{lll}
 \text{first branch:}&u\to u_*<\infty,&v\to+\infty,\\
 \text{second branch:}&v\to v_*<\infty,&u\to+\infty,\\
 \text{corner:}&u\to+\infty,&v\to+\infty.
 \end{array}
\end{equation}
These alternatives describe limiting coordinate regimes; they do not
adjoin boundary points or identify $\mathcal P_d$ with
$\mathcal M_d$.
\end{proposition}

\begin{proof}
The comparison construction is
\cite[\CompanionRef{foundation}{cmp:development}]{Foundation}; the absence of finite
exterior runout and the first-runout classification are respectively
\CompanionRef{foundation}{cmp:completeness} and \CompanionRef{foundation}{cmp:classification} of that companion. Their hypotheses are
exactly the ones imposed here: a fixed rotating subextremal center,
the complete signed two-ended bridge, smooth vacuum data with finite
weighted symbol seminorms, and smallness only in $p_{10}$. The
weighted seminorms and asymptotic powers agree with those of our phase
space; equivalent reference atlases can be accommodated by reducing
the radius. There is no requirement that asymptotic charges be fixed,
that a curvature bound hold, or that an extension exist. The cylinder
and wedge are furnished by the source's finite bridge evolution,
two exterior evolutions, and full-cylinder interior continuation.

To make explicit how the cited classification is used, let
$\tau_P$ be the supremum of the times for which the initial connected
geodesic segment stays in $\mathcal P_d$. Then
$\tau_P\le T_d(w)<\infty$. Viewed in $\mathcal P_d$, this is
a maximal future geodesic segment. The cited exterior completeness
proposition rules out its avoiding $\mathcal S$. Both $u$ and $v$
strictly increase on a future timelike wedge track. If they had finite
limits, its position would remain in a compact regular coordinate
range; the temporal-chart argument of
Lemma~\ref{test:interior-endpoint} would give geodesic continuation
there. Hence one or both limits are infinite, exactly as in
Equation~\eqref{test:runout}. If $\tau_P<T_d(w)$, its endpoint
is an ordinary point of $\mathcal M_d$, although it is not a point
of this wedge. The conclusion does not change in that case.
\end{proof}

For the corner argument we record the low geometric estimates from
\cite[\CompanionRef{foundation}{cmp:corner-low-input} and \CompanionRef{foundation}{cmp:null-connections}]{Foundation} by their
content. Put $L_0=\partial_u+b^A\partial_{\theta^A}$ and
$L_1=\partial_v$, and write $y_0=u$, $y_1=v$. For angular vectors
define $\chi_i=\tfrac12\mathcal L_{L_i}\gamma$,
$\ell_i=L_i\log a$, and the angular one-forms $\xi_i$ by
\begin{equation}\label{test:torsion-def}
 \xi_0+\xi_1=d_\theta\log a,
 \qquad
 \xi_0-\xi_1=-a^{-1}\gamma(\partial_vb,\cdot).
\end{equation}
They obey $|\xi_i|_\gamma\le C$, and there are nonnegative
$g_i\in L^1(\mathbb R)$ with
\begin{equation}\label{test:corner-low}
 |\chi_i|_\gamma+|\ell_i+\kappa|
       \le C(q+g_i(y_i)),
 \qquad |b|_{\mathrm{round}}\le C(q+g_0(u)).
\end{equation}
These are the source's interior low estimates on the actual full
wedge, with no higher smallness or curvature assumption. Koszul's
formula gives, for an angular vector $E$ and $j\ne i$,
\begin{equation}\label{test:null-connections}
 \nabla_{L_i}L_i=\ell_iL_i,
 \qquad
 \nabla_E L_i=\chi_i(E)^\sharp+\xi_i(E)L_i,
 \qquad
 \nabla_{L_j}L_i=a\xi_j^\sharp.
\end{equation}
In particular these identities use the shift in $L_0$ rather than
identifying $L_0$ with $\partial_u$.

The variables $Y_i=e^{-\kappa y_i}$ tend to zero when the corresponding
optical coordinate diverges. We will follow corner tracks to the null
sections $u=j$: their remaining optical coordinate $Y_1$ tends to
zero, while the next estimate bounds their momenta conjugate to the
coordinates on those sections. Conservation of canonical volume will exclude an open
family of corner seeds.

\begin{lemma}[Bounds on each individual wedge track]
\label{test:corner-bounds}
For a future unit timelike geodesic beginning at a fixed point of the
wedge, set $Y_i=e^{-\kappa y_i}$ and
$\Xi=\dot\theta-b\dot u$. On its remaining wedge track each
$|\dot Y_i|$ is bounded above and below by positive constants,
and $|\Xi|_\gamma$ is bounded. The constants may depend on its
initial unit tangent. In particular every corner track has one finite
bound for both $\dot Y_i$ and $\Xi$ throughout a sufficiently
late tail.
\end{lemma}

\begin{proof}
This is the part of the corner-ratio lemma
\cite[\CompanionRef{foundation}{cmp:corner-ratios}]{Foundation} that we need; we include its
estimate to display the dependence on the individual track. Write
\[
 U=\dot y_0L_0+\dot y_1L_1+\Xi,
 \qquad p_i=-g(U,L_i)=a\dot y_j>0\quad(j=1-i).
\]
The future unit condition and Equation~\eqref{test:null-connections}
give
\begin{align}
 1+|\Xi|_\gamma^2&=2a\dot y_0\dot y_1,
       \label{test:normalization}\\
 \frac{\dot p_i}{p_i}
   &=\ell_i\dot y_i+(\xi_i-\xi_j)(\Xi)
                         -\frac{\chi_i(\Xi,\Xi)}{p_i}.
       \label{test:momentum}
\end{align}
For example, the last equation follows by differentiating
$-g(U,L_i)$, using $\nabla_UU=0$ and the three terms in
$\nabla_UL_i$ from Equation~\eqref{test:null-connections}.
Since $|\Xi|_\gamma^2/p_i\le2\dot y_i$, it follows that
\begin{equation}\label{test:logarithmic-variation}
 \left|\frac{d}{d\tau}\log(e^{\kappa y_i}p_i)\right|
  \le C(q+g_i(y_i))\dot y_i+C|\Xi|_\gamma,
 \qquad
 e^{\kappa y_i}p_i=\frac{q_0}{\kappa}\mathcal A|\dot Y_j|.
\end{equation}
The positive functions $Y_i$ decrease. On every initial part of the
remaining track, Equation~\eqref{test:normalization} implies
\[
 \int|\Xi|_\gamma\,d\tau
 \le C\left(\int|\dot Y_0|\,d\tau
                   \int|\dot Y_1|\,d\tau\right)^{1/2}
 \le C\sqrt{Y_0^{\mathrm{in}}Y_1^{\mathrm{in}}}.
\]
Moreover,
\[
 \int q\dot y_i\,d\tau
    \le\frac{q_0}{\kappa}Y_0^{\mathrm{in}}Y_1^{\mathrm{in}},
 \qquad
 \int g_i(y_i)\dot y_i\,d\tau
    \le\int_{y_i^{\mathrm{in}}}^{\infty}g_i(s)\,ds.
\]
Thus the logarithmic variation in
Equation~\eqref{test:logarithmic-variation} is bounded independently
of the terminal point of this track portion. Boundedness of
$\mathcal A^{\pm1}$ gives positive upper and lower bounds for each
$|\dot Y_i|$, proportional to its initial value. Finally
\[
 1+|\Xi|_\gamma^2
   =\frac{2q_0\mathcal A}{\kappa^2}
                 |\dot Y_0|\,|\dot Y_1|
\]
bounds the screen velocity. These constants depend on the two initial
$Y$-speeds; no bound uniform over all unit timelike tangents has been
asserted.
\end{proof}

\subsection{Null sections and exclusion of corner seeds}

\begin{proposition}[Corner seeds have measure zero]
\label{test:corner-null}
For a fixed $d\in\mathcal U$, the seeds whose first runout from
$\mathcal P_d$ is a corner are contained in a set of measure zero
for every smooth positive measure on $T\Sigma$.
\end{proposition}

\begin{proof}
Use $Y_i=e^{-\kappa y_i}$ as above, and consider the interior null
sections
\[
 \mathcal N_j=\{u=j,\ u+v>0\},\qquad j=1,2,\ldots.
\]
Every corner track hits all sufficiently large $\mathcal N_j$
transversely before first runout. The tangential covariant momenta
on these sections are
\begin{equation}\label{test:section-momenta}
 \pi_{Y_1}=-\frac{a\dot Y_0}{\kappa^2Y_0Y_1}
       =-\frac{q_0\mathcal A}{\kappa^2}\dot Y_0,
 \qquad P_A=\gamma_{AB}\Xi^B.
\end{equation}
The shift contributes to the $Y_0$ covector component, but not to
$\pi_{Y_1}$, because $\partial_{Y_1}$ is a multiple of $L_1$.
Lemma~\ref{test:corner-bounds} and the uniform sphere-metric
comparison show that, along each corner track, these momenta satisfy
one finite bound on its entire sufficiently late tail. Also
$Y_1\to0$ along that tail.

Here is the exact volume statement used to transfer these bounds.
On $T^*\mathcal P_d$ take the canonical symplectic form
$\omega=d\pi_\alpha\wedge dx^\alpha$, with the corresponding
Hamiltonian-flow sign convention, and set
$\mathcal H=\tfrac12g^{-1}(\pi,\pi)$. Restrict to the future
unit mass shell $\mathcal H=-\tfrac12$. Above a fixed
$\mathcal N_j$ the six variables
\[
 (Y_1,\theta^1,\theta^2,\pi_{Y_1},P_1,P_2)
\]
give coordinates on the transverse section of this shell. To check
the missing momentum explicitly, let
$\widetilde a=a/(\kappa^2Y_0Y_1)>0$ and
$\widetilde b=-b/(\kappa Y_0)$. Then
\[
 g=-2\widetilde a\,dY_0dY_1+
 \gamma_{AB}(d\theta^A-\widetilde b^A dY_0)
             (d\theta^B-\widetilde b^B dY_0),
\]
and the shell equation is
\begin{equation}\label{test:mass-shell-solve}
 -\widetilde a^{-1}
       (\pi_{Y_0}+\widetilde b^AP_A)\pi_{Y_1}
       +\tfrac12\gamma^{AB}P_AP_B=-\tfrac12.
\end{equation}
Future timelikeness gives $\dot Y_0<0$ and hence
$\pi_{Y_1}=-\widetilde a\dot Y_0>0$. Thus
Equation~\eqref{test:mass-shell-solve} determines
$\pi_{Y_0}$ uniquely. Pulling back $\omega$ to this section,
where $dY_0=0$, gives the nondegenerate form
\begin{equation}\label{test:section-form}
 \omega_j=d\pi_{Y_1}\wedge dY_1
              +\sum_{A=1}^2dP_A\wedge d\theta^A.
\end{equation}
Its absolute symplectic volume is ordinary canonical coordinate
volume; there is no factor involving $a$ or $Y_0$.

Let $D_j\subset T\Sigma$ consist of seeds whose initial connected
track in $\mathcal P_d$ hits $\mathcal N_j$ at a finite time.
This is open: a given compact geodesic segment stays in the open
comparison spacetime, and its transverse hit persists by smooth flow
dependence and the implicit function theorem. The hitting map
$F_j:D_j\to\{\mathcal H=-\tfrac12\}\big|_{\mathcal N_j}$
is smooth. It is injective. A wedge trajectory has strictly
decreasing $Y_0$ and so hits that section at most once; a maximal
trajectory in the globally hyperbolic $\mathcal P_d$ meets its
Cauchy hypersurface $\Sigma$ once. If two hitting states were equal,
uniqueness of the smooth Hamiltonian flow would therefore give the
same initial seed.

On the shell above $\Sigma$, restriction of $\omega$ is the
canonical form on $T^*\Sigma$. Indeed the tangential initial
covector is $\pi|_{T\Sigma}=h(w,\cdot)$, so the map from seeds
is a smooth fiberwise isomorphism onto $T^*\Sigma$. Denote the
pullback of its positive canonical volume by $\mu$; this is a
smooth positive measure on $T\Sigma$.

The hitting map preserves the section symplectic forms and hence
their volumes, including its variable hitting time. In detail,
Hamiltonian flow preserves $\omega$, and its generating vector
$X_{\mathcal H}$ belongs to the kernel of $\omega$ restricted
to the energy shell, since its contraction is, up to sign,
$d\mathcal H$. For a hitting time $\tau(w)$,
\[
 dF_j(V)=d\Phi_{\tau(w)}(V)
               +X_{\mathcal H}(F_j(w))\,d\tau(V).
\]
All terms involving $d\tau$ therefore vanish when evaluated by
$\omega$ on shell-tangent vectors. This proves preservation of
the section forms. Their nondegeneracy also shows that $F_j$ is a
local diffeomorphism. Its injectivity makes it a diffeomorphism onto
its open image, so volume comparison for measurable subsets follows
from change of variables, without any multiplicity factor.

Fix a smooth round norm on angular covectors. For $r>0$ and an
integer $N\ge1$, the region on the target section defined by
\[
 0<Y_1<r,\qquad |\pi_{Y_1}|\le N,
                    \qquad |P|_{\mathrm{round}}\le N
\]
has canonical volume at most $C_Nr$, uniformly in $j$. A finite
angular coordinate cover with precompact chart cores proves this
directly from Equation~\eqref{test:section-form}; covector norms and
coordinate components are uniformly comparable on each core.
Restricting to the actual wedge $u+v>0$ or to the future shell only
reduces that volume. Let $S_{j,r,N}$ be its inverse image under
$F_j$, including only the finite hits in $D_j$. These are measurable
and satisfy
\begin{equation}\label{test:strip-volume}
 \mu(S_{j,r,N})\le C_Nr.
\end{equation}
Fatou's inequality applied to their indicator functions gives
\begin{equation}\label{test:fatou}
 \mu\!\left(\liminf_{j\to\infty}S_{j,r,N}\right)
 \le\liminf_{j\to\infty}\mu(S_{j,r,N})\le C_Nr.
\end{equation}
This use of Fatou does not require $\mu(T\Sigma)$ to be finite.

Every corner seed is contained in
\begin{equation}\label{test:countable-corner-cover}
 \bigcup_{N=1}^{\infty}\ \bigcap_{s=1}^{\infty}
                   \liminf_{j\to\infty}S_{j,1/s,N}.
\end{equation}
Indeed first choose one integer $N$ larger than the eventual momentum
bound of its track. For each $s$ all sufficiently late hits then have
$Y_1<1/s$ and that same momentum bound. For a fixed $N$ the
intersection in Equation~\eqref{test:countable-corner-cover} has
measure at most $C_N/s$ for every $s$, and is therefore null.
The countable union is null. All smooth positive measures on
$T\Sigma$ have the same null sets, proving the assertion. In
particular no uniform momentum bound over the family of corner seeds
has been inserted into the proof.
\end{proof}

\begin{corollary}[A branch seed for every passing test]
\label{test:branch-seed}
If $d\in\mathcal A(O,B)$, then there is $w\in O$ whose first
runout from $\mathcal P_d$ is one of the two branches in
Equation~\eqref{test:runout}. After exchanging the optical labels
and preparing the angular coordinates as in
\cite[\CompanionRef{foundation}{in:label-swap}]{Foundation},
the branch can be written
$u\to u_*<\infty$, $v\to\infty$. The seed is still the
same initial seed in $O$.
\end{corollary}

\begin{proof}
The passing test gives finite full-development duration for every
$w\in O$, so Proposition~\ref{test:comparison} applies to all of
them. The nonempty open set $O$ has positive $\mu$-measure and
cannot be contained in the null set of
Proposition~\ref{test:corner-null}. Choose a noncorner seed. The
source's exchange lemma changes the coordinate preparation on the
same cylinder and product spacetime and retains its estimates; it
changes neither the geodesic nor its initial seed.
\end{proof}

The selected branch supplies an observer for the finite perturbation
argument. At this stage the only bounds imposed on its full future
track are the lifetime and transport bounds in
Definition~\ref{test:definition}; no curvature bound or regular
ambient endpoint of this observer is required.

The remaining task can now be stated precisely. For every passing datum
$d\in\mathcal A(O,B)$, finite $m\ge10$, and $r>0$, we must construct
$d'\in\mathcal U$ such that
\begin{equation}\label{test:density-target}
 p_m(d'-d)<r,\qquad d'\notin\mathcal A(O,B).
\end{equation}
By Lemma~\ref{test:closed}, this makes each test closed with empty
relative interior. The Baire property and Proposition~\ref{test:exit-cover}
then exclude future extensions on the countable intersection of their
open dense complements. Sections~\ref{sec:background}--\ref{holo:section}
construct $d'$ by a finite experiment near a late point of the selected
branch; Section~\ref{sec:completion} completes this assembly.
